\documentclass{article}
\usepackage{amsmath}
\title{Rotation Convention and Formula Derivations}
\author{Giacomo Accorto}
\date{\today}

\begin{document}
\maketitle

\section{Reference Frame and Angle Conventions}

We define a right-handed, vehicle-fixed, coordinate frame attached to the moving plane:
\begin{itemize}
  \item $x_v$: forward direction of the vehicle,
  \item $y_v$: lateral axis (pointing to the left),
  \item $z_v$: normal to the rover plane, pointing upward toward the GNSS module.
\end{itemize}

The GNSS module is located at a fixed height $h = 1500\ \text{mm}$ along the $z_v$ axis
from the rover plane.

Roll and pitch are defined according to the task statement:
\begin{itemize}
  \item \textbf{Roll} ($\phi$): rotation about $x_v$, positive when the right side of
  the vehicle is lower than the left. Assume
  $-\frac{\pi}{2} < \phi < \frac{\pi}{2}$.
  \item \textbf{Pitch} ($\theta$): rotation about $y_v$, positive when the front of the
  vehicle is lower than the rear. Assume
  $-\frac{\pi}{2} < \theta < \frac{\pi}{2}$.
\end{itemize}
This implies that all rotations follow the right-hand rule.

\subsection{Rotation Convention}

The rover orientation is constructed using intrinsic (body-fixed) rotations:
\begin{itemize}
  \item first a roll rotation about $x_v$,
  \item then a pitch rotation about the updated $y_v$ axis.
\end{itemize}

This rotation order is assumed, as it is not explicitly specified in the task statement.
It is interesting to notice that, for small angles (smooth movement), the rotation order
becomes a second-order effect ($R_\theta R_\phi \approx R_\phi R_\theta$ ).

\section{Projection of the GNSS Module on the Vehicle Plane}

\subsection{Pitch Rotation}

Define a frame of reference v at the base of the post.
The GNSS module position when the vehicle is flat is simply:
\begin{equation}
\vec{h}_v = (0, 0, h)
\end{equation}
In this frame of reference, a roll rotation $\phi$ around the x-axis is given by the operator
\begin{equation}
  R_x(\phi)=
\begin{bmatrix}
1 & 0 & 0 \\
0 & \cos\phi & -\sin\phi \\
0 & \sin\phi & \cos\phi
\end{bmatrix}
\end{equation}
which, applied to the vector representing the GNSS module of the vehicle, gives
\begin{equation}
R_x(\phi)\vec{h}_v=(0, -h\sin\phi, h\cos\phi)
\end{equation}
This correctly shows that positive roll tilts the z-axis to smaller y, meaning the right side is lower than the left one.

\subsection{Pitch Rotation}

The pitch rotation can be represented instead by the operator:
\begin{equation}
R_y(\theta)=
\begin{bmatrix}
\cos\theta & 0 & \sin\theta \\
0 & 1 & 0 \\
-\sin\theta & 0 & \cos\theta
\end{bmatrix}
\end{equation}
which, applied to the GNSS post vector gives
\begin{equation}
R_y(\theta)\vec{h}_v=(h\sin\theta, 0, h\cos\theta)
\end{equation}
This correctly shows that positive pitch tilts the z-axis to larger x, meaning the front part of the vehicle is being lower
than the rear part.

\subsection{Combined Rotation and Projection}

Applying both operations (roll, then pitch) to the post vector gives:
\begin{equation}
  R_y(\theta)R_x(\phi)\vec{h}_v=(h\sin\theta\cos\phi,\,-h\sin\phi,\,h\cos\theta\cos\phi)
\end{equation}

The data gives us the coordinate of the GNSS (G) module in the word frame of reference: $(x_G, y_G, z_G)$
To find the base of the post, we simply consider that the coordinate $(x_G, y_G, z_G) = (x_P, y_P, z_P) + (h\sin\theta\cos\phi,\,-h\sin\phi,\,h\cos\theta\cos\phi)$
From this formula, we read out the projection of the post on the vehicle:
\begin{equation}
  (x_P, y_P) =  (x_G- h\sin\theta\cos\phi, y_G+ h\sin\phi)
\end{equation}

\section{Velocity and Heading Estimation}

Once the GNSS module projection on the moving plane has been computed, the vehicle
heading can be calculate from the direction of motion (velocity) in the world frame. Let
$(x_P(t), y_P(t))$ be the projected point trajectory. Then the planar velocity is
\begin{equation}
  \vec{v}(t) =
  \begin{pmatrix}
    v_x(t) \\
    v_y(t)
  \end{pmatrix}
\end{equation}
where $v_x(t)$ and $v_y(t)$ are approximated numerically from neighboring data.

Assuming the vehicle moves smoothly forward, the heading angle is taken as the angle
of this planar velocity vector:
\begin{equation}
  \psi(t) = \operatorname{atan2}\bigl(v_y(t), v_x(t)\bigr)
\end{equation}
where $\operatorname{atan2}(y,x)$ is the quadrant-aware $\arctan$, the angle $\psi \in (-\pi,\pi]$ such that
\[
\cos\psi = \frac{x}{\sqrt{x^2+y^2}},
\qquad
\sin\psi = \frac{y}{\sqrt{x^2+y^2}}.
\]
This gives the heading in the world frame, measured from the positive $x$-axis.


\end{document}
